\begin{outline}
What is different between designing an authoritative storage system and designing a permissionless one?

The storage systems do essentially same things to ensure durability.
Store with redundancy.
Detect redundancy loss in time.
Recover expected redundancy level in time.
Durability proved by showing the recovery is always faster than the series of loss arriving at an independently random interval.

The difference lies in studying the correlations of the failures.
There's no way to study the correlations in a permissionless system in a reliable and trusted way.
Correlations in the wild are multidimensional.
In datacenters storage nodes are carefully deployed to reduce the dimensions.
This is impossible in permissionless systems.

We must consider \emph{all} possible correlated failures in a permissionless system.
Every combination of nodes may fail together, while they may still fail by themselves.
Different combinations may fail with different probabilities.
But generally nothing is impossible, since we can assert nothing about the participant nodes.

Define failure generator.

Define fault.

Define fault tolerance.

The failure generator may assign certain failure domains with particular high failure rates out of two reasons.
The static reason is the inherent failure patterns of the correlated nodes.
The dynamical/adversarial reason is that attackers may exploit system's assumptions on failure domains as vulnerabilities.
The ability of attackers manipulating failure generator varies, the detailed security analysis is out of topic.
The important piece is that the effect of two factors are indistinguishable.
We cannot rely on things like ``it is an honest node but failed out of nature disaster so it is easier to be tolerated''.

We don't have an explicit model for adversarial/arbitrary faults in this work.
In our proposed protocol there's no coordination, so it (mostly) does not rely on \eg whether the interacted nodes behave correctly.
Those faulty behaviors are also covered by correlated failures.
\end{outline}